\chapter{Vergleichsinterfaces: \texttt{std::map} und \texttt{std::vector}}\label{app:interfaces}

% Anhang F: vollständige C++23-Schnittstellen der beiden Vergleichshüllen (std::map, std::vector) zu
% Abschnitt sec:measurement-system; Semantik, Komplexität und Fassung nach ISO/IEC 14882:2024
% (iso_cpp), Sekundärquelle cppreference (cppreference_map, cppreference_vector).
% Dagger = Kern-Operationen, die das Konformitäts-Gatter der jeweiligen Gattung vor jeder Messung
% prüft (Map-Gattung: std::map-Orakel, Blöcke RF1..RF14, davon RF14 der Opt-in-Block wide_keys;
% Sequenz-Gattung: std::deque-Orakel, Blöcke QRF1..QRF8, davon QRF8 der als übersprungen
% protokollierte Block push_front/pop_front).
% Entwickler-Dokument der Cache-Engine: docs/ENTWICKLER-FUNCTION-HANDLE-HOPS.md (Sprungpfad je Funktion).
Dieser Anhang führt --- ergänzend zu Abschnitt~\ref{sec:measurement-system} --- die vollständige
Auflistung aller Standard-Interface-Funktionen der beiden Vergleichshüllen \texttt{std::map}
(geordnet-assoziativ) und \texttt{std::vector} (Sequenz) nach dem Stand von C++23 mit ihrer
offiziell erwarteten Semantik und ihrer Komplexitäts-Garantie gemäß dem
ISO-C++-Standard~\cite{iso_cpp}. Er dient zugleich als zentrale Referenz für die Weiterentwicklung
der Cache-Engine: Da die Suchalgorithmus-Hülle ihr Innenleben zur Übersetzungszeit austauscht ---
je Achsen-Wahl entsteht ein eigenes Kompilat ---, während zur Laufzeit nur die Parameter der
dynamischen Unter-Achsen nachgestellt werden (die registrierten dynamischen Dimensionen des
Mess-Realms, Abschnitt~\ref{sec:measurement-system}), wird je Standard-Funktion
erklärungsbedürftig, welches Verhalten der äußere Vertrag garantiert.

\emph{Variadische Abbildung.} Die hybride Such-Engine des Prüflings PRT-ART bindet die
Parameterzahl an die Hülle: Ein Template-Parameter fordert alle Funktionsinterfaces eines
\texttt{std::vector}, zwei die eines \texttt{std::map} und $N>2$ die einer
\texttt{std::map<K,\,std::tuple<\ldots>>} (Abschnitt~\ref{sec:impl-architecture}). Die Cache-Engine
leitet aus derselben Regel die Typen ab --- bei einem Parameter wird der Schlüssel implizit als
monoton gezählter 64-Bit-Wert vergeben, bei $N>2$ wird der Wert zum Tupel ---; ihre beiden
Vertrags-Hüllen sind als getrennte Gattungs-Adapter gebaut: die Map-Gattung über den funktionalen
Antrieb der Modulgrenze mit den fünf Operationen \texttt{tier\_insert}, \texttt{tier\_lookup},
\texttt{tier\_erase}, \texttt{tier\_clear} und \texttt{tier\_size}, die Sequenz-Gattung über den
Sequenz-Antrieb mit \texttt{tier\_push\_back}, \texttt{tier\_at}, \texttt{tier\_size} und
\texttt{tier\_clear}; ein Adapter quer zur Gattung ist im Typsystem ausgeschlossen. Der Fall $N>2$
ist Typ-Ableitung ohne eigenen Adapter; die Modulgrenze führt Schlüssel und Wert durchgängig als
64-Bit-Werte. Die geteilten Funktionen (Iteration, \texttt{size}/\texttt{empty}/\texttt{clear},
\texttt{swap}, \texttt{get\_allocator}) tragen in beiden Hüllen denselben Namen und dieselbe Rolle;
über die Modulgrenze reichen beide Hüllen davon \texttt{size} und \texttt{clear}, \texttt{empty} ist
host-seitig abgeleitet, Iteration und \texttt{swap} queren die Grenze nicht, und
\texttt{get\_allocator} ist auf der Map-Seite ein optionales Proxy-Sub-Interface außerhalb des
Orakels (Abschnitt~\ref{asec:if-decomp}).

\emph{Zusammengesetzte Funktionen.} Einige Funktionen ergeben sich \enquote{unter der Haube} als Komposition
primitiverer Operationen --- etwa \texttt{operator[]} als \texttt{lower\_bound} mit anschließendem
Default-Einfügen bei fehlendem Schlüssel oder \texttt{at} als \texttt{find} mit Wurf einer
\texttt{std::out\_of\_range}-Ausnahme bei Fehlen ---; solche Fälle sind in der Spalte \enquote{Anmerkung}
gekennzeichnet. Die Spalte \enquote{seit} nennt die einführende Standard-Fassung. Drei Begriffe der
Tafeln stammen aus der jüngeren Standard-Entwicklung: der Node-Handle\footnote{Seit C++17 kann ein
Knoten eines assoziativen Containers mit \texttt{extract} ohne Kopie oder Move als Node-Handle gelöst
und mit \texttt{insert} bzw.\ \texttt{merge} in einen anderen Container gespleißt
werden~\cite{talbot2016splicing}.}, die Deduktions-Leitfäden (CTAD)\footnote{\emph{Class template
argument deduction}, seit C++17: Die Template-Argumente eines Klassentemplates werden aus den
Konstruktor-Argumenten abgeleitet, über implizit aus den Konstruktoren gebildete oder explizit
angegebene Leitfäden~\cite{spertus2016ctad}.} und der transparente Schlüssel
(Abschnitt~\ref{asec:if-map}). Die exakten cache-engine-internen Funktions-Sprünge
(\emph{function-handle-hops}\footnote{Hausbegriff der Cache-Engine für die Folge der internen
Funktions-Sprünge (Handles) vom Eintritt an der Modulgrenze über die Achsen-Organe bis zum Ergebnis;
das Entwickler-Dokument verzeichnet sie je Standard-Funktion und markiert nicht verdrahtete
Funktionen als solche.}) dokumentiert das Entwickler-Dokument
\texttt{docs/\allowbreak{}ENTWICKLER-\allowbreak{}FUNCTION-\allowbreak{}HANDLE-\allowbreak{}HOPS.md}
der Cache-Engine; es verzeichnet je Standard-Funktion den Sprungpfad und den Status an der
Modulgrenze.

\emph{Konformitäts-Teilmenge.} Das Konformitäts-Gatter (Abschnitt~\ref{sec:measurement-system})
prüft vor jeder Messung je Gattung eine Kern-Teilmenge der Hüllen-Operationen gegen ein Orakel der
Standardbibliothek --- \texttt{std::map} für die Map-Gattung, \texttt{std::deque} für die
Sequenz-Gattung ---, damit nur Vergleichbares verglichen wird; die Teilmenge ist in beiden Tafeln
mit $\dagger$ gekennzeichnet.

\section{\texttt{std::map} (geordnet-assoziativ)}\label{asec:if-map}
Tabelle~\ref{tab:if-map} listet die vollständige Schnittstelle von
\texttt{std::map<Key,\,T,\,Compare,\,Allocator>} nach C++23. Ein $\dagger$ markiert die
Kern-Operationen, die das Konformitäts-Gatter (Abschnitt~\ref{sec:measurement-system}) vor jeder
Messung prüft, damit nur Vergleichbares verglichen wird. Komplexitäten beziehen sich auf
$n=\texttt{size()}$; bei Bereichs-Operationen ist $N$ die Zahl der Eingabe-Elemente; Fassungs-Marker
und Komplexitäten folgen dem Normtext~\cite{iso_cpp} und sind an der
Online-Referenz~\cite{cppreference_map} gegengelesen. Die siebzehn markierten Operationen sind
genau die Operationsmenge der zwölf Kern-Prüfblöcke des Gatters: Keine markierte Operation bleibt
ungeprüft, keine geprüfte Operation ist unmarkiert (\texttt{size} wird in den Blöcken mitgeprüft,
gehört aber zum Zufalls-Kernblock und ist deshalb nicht markiert). Vier der markierten Operationen
entsprechen unmittelbar Operationen der Modulgrenze (\texttt{insert\_or\_assign} dem
Einfüge-Antrieb, \texttt{find} der Punkt-Suche, \texttt{erase} und \texttt{clear} ihren Antrieben);
\texttt{at}, \texttt{operator[]}, \texttt{count}, \texttt{contains}, \texttt{empty} sowie
\texttt{insert}, \texttt{emplace} und \texttt{try\_emplace} werden host-seitig komponiert ---
\texttt{empty} aus der Größe, die übrigen aus Punkt-Suche und Einfügen; die drei nicht
überschreibenden Formen laufen über einen Pfad und werden gegen drei verschiedene Orakel-Aufrufe
gestellt; \texttt{begin}/\texttt{end}, \texttt{lower\_bound}, \texttt{upper\_bound} und
\texttt{equal\_range} synthetisiert der Host über die additive Scan-Fähigkeit der Modulgrenze
geordnet, und fehlt sie einer geladenen Binary, protokolliert das Gatter diese Blöcke als
übersprungen, nicht als bestanden; \texttt{key\_comp} wird am Orakel festgestellt (Schlüsselordnung
kleiner-als auf 64-Bit-Werten), weil die Hülle keinen Komparator über die Grenze trägt.

\begin{longtable}{@{}>{\raggedright\arraybackslash}p{3.1cm} >{\raggedright\arraybackslash}p{5.7cm} >{\raggedright\arraybackslash}p{2.7cm} >{\raggedright\arraybackslash}p{1.1cm}@{}}
  \caption{Vollständige \texttt{std::map}-Schnittstelle (C++23)}\label{tab:if-map}\\
  \toprule
  \textbf{Funktion} & \textbf{Offizielle Semantik} & \textbf{Komplexität} & \textbf{seit} \\
  \midrule
\endfirsthead%
  \toprule
  \textbf{Funktion} & \textbf{Offizielle Semantik} & \textbf{Komplexität} & \textbf{seit} \\
  \midrule
\endhead%
  \multicolumn{4}{@{}l}{\textbf{Konstruktion, Zuweisung, Allokator}}\\
  \addlinespace[2pt]
  \texttt{(Konstruktor)} & leer; aus \texttt{Compare}/\texttt{Allocator}; Bereich; Kopie/Move (auch allokator-erweitert); \texttt{initializer\_list}; ab C++23 \texttt{from\_range} & $O(1)$ bis $O(N\log N)$ & C++98 \\
  \texttt{(Destruktor)} & zerstört alle Elemente und Knoten & $O(n)$ & C++98 \\
  \texttt{operator=} & Kopier-/Move-/\texttt{initializer\_list}-Zuweisung mit Allokator-Propagation & $O(n)$ & C++98 \\
  \texttt{get\_allocator} & Kopie des zugeordneten Allokators & $O(1)$ & C++98 \\
  \addlinespace[3pt]
  \multicolumn{4}{@{}l}{\textbf{Elementzugriff}}\\
  \addlinespace[2pt]
  \texttt{at}\,$\dagger$ & Referenz auf den Mapped-Wert; wirft \texttt{std::out\_of\_range} bei Fehlen & $O(\log n)$ & C++11 \\
  \texttt{operator[]}\,$\dagger$ & Referenz; fügt bei Fehlen wertinitialisiert ein (kein \texttt{std::out\_of\_range}; das Einfügen kann \texttt{std::bad\_alloc} oder Konstruktor-Ausnahmen werfen) & $O(\log n)$ & C++98 \\
  \addlinespace[3pt]
  \multicolumn{4}{@{}l}{\textbf{Iteratoren}}\\
  \addlinespace[2pt]
  \texttt{begin}/\texttt{end}\,$\dagger$ & Iterator auf erstes Element bzw.\ Ende, in Schlüsselordnung (\texttt{cbegin}/\texttt{cend} ab C++11) & $O(1)$ & C++98 \\
  \texttt{rbegin}/\texttt{rend} & Rückwärts-Iteratoren (\texttt{crbegin}/\texttt{crend} ab C++11) & $O(1)$ & C++98 \\
  \addlinespace[3pt]
  \multicolumn{4}{@{}l}{\textbf{Kapazität}}\\
  \addlinespace[2pt]
  \texttt{empty}\,$\dagger$ & leer-Test ($\texttt{begin()==end()}$); \texttt{[[nodiscard]]} ab C++20 & $O(1)$ & C++98 \\
  \texttt{size} & Anzahl der Elemente & $O(1)$ & C++98 \\
  \texttt{max\_size} & theoretische Maximalgröße & $O(1)$ & C++98 \\
  \addlinespace[3pt]
  \multicolumn{4}{@{}l}{\textbf{Modifikatoren}}\\
  \addlinespace[2pt]
  \texttt{clear}\,$\dagger$ & alle Elemente entfernen & $O(n)$ & C++98 \\
  \texttt{insert}\,$\dagger$ & Element(e) ohne Überschreiben (Wert, \texttt{P\&\&}, Hint, Bereich, \texttt{initializer\_list}, Node-Handle ab C++17) & $O(\log n)$; Hint amort.\ $O(1)$; Bereich $O(N\log(n{+}N))$ & C++98 \\
  \texttt{insert\_range} & Eingabe-Bereich einfügen, keine Überschreibung & $O(N\log(n{+}N))$ & C++23 \\
  \texttt{insert\_\allowbreak{}or\_\allowbreak{}assign}\,$\dagger$ & einfügen oder vorhandenen Wert zuweisen; \texttt{bool} meldet Einfügung & $O(\log n)$ & C++17 \\
  \texttt{emplace}\,$\dagger$ & in-place konstruieren, nur falls Schlüssel fehlt (sonst verworfen) & $O(\log n)$ & C++11 \\
  \texttt{emplace\_hint} & wie \texttt{emplace}, positionsnah zum Hint & $O(\log n)$; amort.\ $O(1)$ & C++11 \\
  \texttt{try\_emplace}\,$\dagger$ & wie \texttt{emplace}, rührt die Argumente bei vorhandenem Schlüssel nicht an & $O(\log n)$ & C++17 \\
  \texttt{erase}\,$\dagger$ & nach Iterator/Bereich/Schlüssel (Schlüssel: 0 oder 1 entfernt) & Iter.\ amort.\ $O(1)$; Key $O(\log n)$ & C++98 \\
  \texttt{swap} & Inhalt und Comparator tauschen; Iteratoren bleiben gültig & $O(1)$ & C++98 \\
  \texttt{extract} & Element als Node-Handle entnehmen (ohne Kopie/Move) & Iter.\ amort.\ $O(1)$; Key $O(\log n)$ & C++17 \\
  \texttt{merge} & Knoten aus Quell-Map splicen; vorhandene Schlüssel bleiben dort & $O(N\log(n{+}N))$ & C++17 \\
  \addlinespace[3pt]
  \multicolumn{4}{@{}l}{\textbf{Suche}}\\
  \addlinespace[2pt]
  \texttt{find}\,$\dagger$ & Iterator auf den Schlüssel oder \texttt{end()} & $O(\log n)$ & C++98 \\
  \texttt{count}\,$\dagger$ & Anzahl (0 oder 1) & $O(\log n)$ & C++98 \\
  \texttt{contains}\,$\dagger$ & Mitgliedschaft (\texttt{bool}) & $O(\log n)$ & C++20 \\
  \texttt{lower\_bound}\,$\dagger$ & erster Schlüssel $\ge$ \texttt{key} & $O(\log n)$ & C++98 \\
  \texttt{upper\_bound}\,$\dagger$ & erster Schlüssel $>$ \texttt{key} & $O(\log n)$ & C++98 \\
  \texttt{equal\_range}\,$\dagger$ & Paar \texttt{[lower\_bound,\allowbreak\,upper\_bound)} & $O(\log n)$ & C++98 \\ % chktex 9
  \addlinespace[3pt]
  \multicolumn{4}{@{}l}{\textbf{Beobachter}}\\
  \addlinespace[2pt]
  \texttt{key\_comp}\,$\dagger$ & Kopie des Schlüssel-Comparators & $O(1)$ & C++98 \\
  \texttt{value\_comp} & \texttt{value\_compare} (vergleicht Paare über den Schlüssel) & $O(1)$ & C++98 \\
  \addlinespace[3pt]
  \multicolumn{4}{@{}l}{\textbf{Nicht-Member-Funktionen}}\\
  \addlinespace[2pt]
  \texttt{operator==} & gleiche Größe und paarweise gleiche Elemente & $O(n)$ & C++98 \\
  \texttt{operator<=>} & lexikografischer Drei-Wege-Vergleich; synthetisiert \texttt{<\,<=\,>\,>=} & $O(n)$ & C++20 \\
  \texttt{std::swap} & Spezialisierung; ruft \texttt{lhs.swap(rhs)} & $O(1)$ & C++98 \\
  \texttt{std::erase\_if} & entfernt alle Elemente mit \texttt{pred==true}; gibt die Anzahl zurück & $O(n)$ & C++20 \\
  Deduktions-Leitfäden & CTAD aus Bereich/\allowbreak\texttt{initializer\_list}/\allowbreak\texttt{from\_range} (\texttt{from\_range} ab C++23) & --- (Compile-Zeit) & C++17 \\
  \bottomrule
\end{longtable}

\noindent Die vor C++20 eigenständigen Vergleichsoperatoren \texttt{<\,<=\,>\,>=} werden seit C++20 aus
\texttt{operator<=>} synthetisiert, \texttt{!=} aus \texttt{operator==} umgeschrieben. Die
heterogenen Überladungen von \texttt{operator[]}, \texttt{at}, \texttt{try\_emplace} und
\texttt{insert\_or\_assign} (transparenter Schlüssel\footnote{Heterogene Suche seit C++14: Trägt der
Komparator den Typ \texttt{is\_transparent}, nehmen die Suchfunktionen Argumente eines anderen Typs
als \texttt{key\_type} an, ohne ein temporäres Schlüsselobjekt zu
bilden~\cite{wakely2013heterogeneous}; C++23 dehnt dies auf \texttt{erase}/\texttt{extract}
aus~\cite{boyarinov2021erasure}, C++26 auf \texttt{operator[]}, \texttt{at}, \texttt{try\_emplace}
und \texttt{insert\_or\_assign}~\cite{boyarinov2023heterogeneous}.}) gehören erst zu C++26 und
bleiben hier ausgeklammert; die C++23-Neuzugänge sind der \texttt{from\_range}-Konstruktor,
\texttt{insert\_range}, die transparenten Schlüssel-Überladungen von
\texttt{erase}/\texttt{extract} sowie die zugehörigen
\texttt{from\_range}-Deduktions-Leitfäden~\cite{jabot2022ranges,boyarinov2021erasure}. Auch die
\texttt{constexpr}-Fähigkeit von \texttt{std::map} gehört erst zu C++26 und bleibt
ausgeklammert~\cite{cppreference_map}.

\section{\texttt{std::vector} (Sequenz)}\label{asec:if-vector}
Tabelle~\ref{tab:if-vector} listet die vollständige Schnittstelle von
\texttt{std::vector<T,\,Allocator>} nach C++23 ($n=\texttt{size()}$, bei Bereichs-Operationen $N$
die Zahl der Eingabe-Elemente; Fassungs-Marker und Komplexitäten folgen dem
Normtext~\cite{iso_cpp} und sind an der Online-Referenz~\cite{cppreference_vector} gegengelesen).
Sämtliche Member sind seit C++20 \texttt{constexpr}\footnote{Seit C++20 sind alle Elementfunktionen
von \texttt{std::vector} \texttt{constexpr}; ein \texttt{constexpr}-Objekt selbst scheitert in der
Regel am dynamischen Speicher~\cite{dionne2019constexprvector}.}. Ein $\dagger$ markiert die
Kern-Operationen, die das Konformitäts-Gatter der Sequenz-Gattung vor jeder Messung gegen ein
\texttt{std::deque}-Orakel prüft (Anhängen, Index-Zugriff, Größe, Leeren; für diese vier
Operationen ist die \texttt{deque}-Semantik mit der des \texttt{std::vector} deckungsgleich): Die
eigene Sequenz-Hülle bietet genau diese vier Funktionen über die Modulgrenze an, der Index-Zugriff
ist dort geprüft statt undefiniert (Status-Rückgabe: Treffer-Flag und Ausgabewert an der
Modulgrenze, \texttt{std::optional} im Organ). Alle übrigen Funktionen der Tafel ---
Kapazitäts-Steuerung, Positions- und Bereichs-Operationen,
\texttt{front}/\texttt{back}/\texttt{pop\_back}, \texttt{data}, Iteratoren, \texttt{swap} und die
Nicht-Member --- queren die Modulgrenze nicht; das Wachstum am vorderen Ende
(\texttt{push\_front}), das das \texttt{deque}-Orakel zusätzlich böte, ist als Ausbau der
Modulgrenze geführt und wird im Gatter als übersprungen protokolliert, nicht als bestanden.

\begin{longtable}{@{}>{\raggedright\arraybackslash}p{3.1cm} >{\raggedright\arraybackslash}p{5.7cm} >{\raggedright\arraybackslash}p{2.7cm} >{\raggedright\arraybackslash}p{1.1cm}@{}}
  \caption{Vollständige \texttt{std::vector}-Schnittstelle (C++23)}\label{tab:if-vector}\\
  \toprule
  \textbf{Funktion} & \textbf{Offizielle Semantik} & \textbf{Komplexität} & \textbf{seit} \\
  \midrule
\endfirsthead%
  \toprule
  \textbf{Funktion} & \textbf{Offizielle Semantik} & \textbf{Komplexität} & \textbf{seit} \\
  \midrule
\endhead%
  \multicolumn{4}{@{}l}{\textbf{Konstruktion, Zuweisung, Allokator}}\\
  \addlinespace[2pt]
  \texttt{(Konstruktor)} & leer/\texttt{Allocator}; $n\times$ Default/Wert; Bereich; Kopie/Move (auch allokator-erweitert); \texttt{initializer\_list}; ab C++23 \texttt{from\_range} & $O(1)$ oder $O(n)$ & C++98 \\
  \texttt{(Destruktor)} & zerstört alle Elemente, gibt Speicher frei & $O(n)$ & C++98 \\
  \texttt{operator=} & Kopie/Move/\texttt{initializer\_list} & $O(n)$ & C++98 \\
  \texttt{assign} & Inhalt durch $n\times$ Wert/\allowbreak{}Bereich/\allowbreak{}\texttt{initializer\_\allowbreak{}list} ersetzen & $O(n)$ & C++98 \\
  \texttt{assign\_range} & Inhalt durch Eingabe-Bereich ersetzen & $O(n)$ & C++23 \\
  \texttt{get\_allocator} & Kopie des zugeordneten Allokators & $O(1)$ & C++98 \\
  \addlinespace[3pt]
  \multicolumn{4}{@{}l}{\textbf{Elementzugriff}}\\
  \addlinespace[2pt]
  \texttt{at}\,$\dagger$ & geprüfter Zugriff; wirft \texttt{std::out\_of\_range} & $O(1)$ & C++98 \\
  \texttt{operator[]} & ungeprüfter Zugriff (UB außerhalb des Bereichs) & $O(1)$ & C++98 \\
  \texttt{front}/\texttt{back} & Referenz auf erstes/letztes Element (UB bei leer) & $O(1)$ & C++98 \\
  \texttt{data} & Zeiger auf den zusammenhängenden Speicher (nicht bei \texttt{vector<bool>}) & $O(1)$ & C++11 \\
  \addlinespace[3pt]
  \multicolumn{4}{@{}l}{\textbf{Iteratoren}}\\
  \addlinespace[2pt]
  \texttt{begin}/\texttt{end}, \texttt{rbegin}/\texttt{rend} & Vorwärts-/Rückwärts-Iteratoren in Indexordnung (\texttt{c…}-Varianten ab C++11) & $O(1)$ & C++98 \\
  \addlinespace[3pt]
  \multicolumn{4}{@{}l}{\textbf{Kapazität}}\\
  \addlinespace[2pt]
  \texttt{empty}/\allowbreak\texttt{size}\,$\dagger$/\allowbreak\texttt{max\_size} & leer-Test, Elementanzahl, Maximalgröße & $O(1)$ & C++98 \\
  \texttt{reserve} & Kapazität $\ge n$; realloziert und invalidiert; wirft \texttt{length\_error} bei $n>$\texttt{max\_size()} & $O(n)$ & C++98 \\
  \texttt{capacity} & aktuell allozierte Kapazität & $O(1)$ & C++98 \\
  \texttt{shrink\_to\_fit} & nicht-bindende Reduktion der Kapazität auf \texttt{size()} & $O(n)$ & C++11 \\
  \addlinespace[3pt]
  \multicolumn{4}{@{}l}{\textbf{Modifikatoren}}\\
  \addlinespace[2pt]
  \texttt{clear}\,$\dagger$ & alle Elemente entfernen (Kapazität bleibt) & $O(n)$ & C++98 \\
  \texttt{insert} & vor Position: Wert/$n\times$/Bereich/\allowbreak\texttt{initializer\_list}; verschiebt den Schwanz & $O(n)$; Bereichsformen $O(n{+}N)$ & C++98 \\
  \texttt{insert\_range} & Eingabe-Bereich vor Position einfügen & $O(n{+}N)$ & C++23 \\
  \texttt{emplace} & Element in-place vor Position konstruieren & $O(n)$ & C++11 \\
  \texttt{erase} & Position/Bereich entfernen, Schwanz nachrücken & $O(n)$ & C++98 \\
  \texttt{push\_back}\,$\dagger$ & Element anhängen (Kopie/Move) & amort.\ $O(1)$ & C++98 \\
  \texttt{emplace\_back} & in-place anhängen; Referenz-Rückgabe ab C++17 & amort.\ $O(1)$ & C++11 \\
  \texttt{append\_range} & Eingabe-Bereich anhängen & amort.\ $O(N)$ & C++23 \\
  \texttt{pop\_back} & letztes Element entfernen (UB bei leer) & $O(1)$ & C++98 \\
  \texttt{resize} & auf $n$ kürzen oder mit Default/Wert auffüllen & $O(n)$ & C++98 \\
  \texttt{swap} & Inhalt und Kapazität mit anderem Vektor tauschen & $O(1)$ & C++98 \\
  \addlinespace[3pt]
  \multicolumn{4}{@{}l}{\textbf{Nicht-Member-Funktionen}}\\
  \addlinespace[2pt]
  \texttt{operator==} & gleiche Größe und elementweise Gleichheit & $O(n)$ & C++98 \\
  \texttt{operator<=>} & lexikografischer Drei-Wege-Vergleich & $O(n)$ & C++20 \\
  \texttt{std::swap} & Spezialisierung; ruft \texttt{lhs.swap(rhs)} & $O(1)$ & C++98 \\
  \texttt{std::erase} & entfernt alle Elemente gleich dem Wert; gibt die Anzahl zurück & $O(n)$ & C++20 \\
  \texttt{std::erase\_if} & entfernt alle Elemente mit \texttt{pred==true}; gibt die Anzahl zurück & $O(n)$ & C++20 \\
  Deduktions-Leitfäden & CTAD aus Bereich/\allowbreak\texttt{from\_range} (\texttt{from\_range} ab C++23) & --- (Compile-Zeit) & C++17 \\
  \bottomrule
\end{longtable}

\noindent Die Spezialisierung \texttt{std::vector<bool>} packt Wahrheitswerte bitweise: Sie besitzt kein
\texttt{data()} und keine Zusammenhangs-Garantie, liefert über \texttt{operator[]}/\texttt{at}/\texttt{front}/%
\texttt{back} einen Proxy-Typ\footnote{Stellvertreter-Objekt: \texttt{vector<bool>::reference}
verhält sich im Ausdruck wie \texttt{bool\&}, bezeichnet aber ein einzelnes Bit statt eines
adressierbaren Objekts, weshalb \texttt{data()} und die Zusammenhangs-Garantie
entfallen~\cite{iso_cpp}.} \texttt{reference} statt \texttt{bool\&} und bietet zusätzlich \texttt{flip()}
(auf Vektor- und auf \texttt{reference}-Ebene), ein statisches
\texttt{swap(reference,\allowbreak\ reference)} sowie eine \texttt{std::hash}-Spezialisierung. Die
Vergleichsoperatoren verhalten sich wie bei \texttt{std::map}.

\section{Zusammengesetzte Funktionen und Konformitäts-Teilmenge}\label{asec:if-decomp}
Mehrere Schnittstellen-Funktionen sind \enquote{unter der Haube} Kompositionen primitiverer Operationen
(Tabelle~\ref{tab:if-decomp}); die Tabelle beschreibt den ISO-Standardvertrag. Die eigene Mess-Hülle
weicht an zwei Stellen bewusst davon ab, weil der ABI-Konformitätsvertrag \texttt{noexcept} verlangt
und eine werfende Hülle als nicht-konform gilt: \texttt{at} liefert statt der Ausnahme eine
Status-Rückgabe (Treffer-Flag und Ausgabewert an der Modulgrenze, \texttt{std::optional} im Organ),
und \texttt{operator[]} ist an der Modulgrenze kein eigener Aufruf, sondern wird host-seitig aus
Punkt-Suche und Default-Einfügen komponiert --- mit der ISO-Semantik (Einfügen bei Fehlen, wirft
nie), aber ohne Referenz-Rückgabe über die binäre Grenze; der Schreibpfad läuft über
\texttt{insert\_or\_assign}. In der Sequenz-Hülle ist \texttt{operator[]} zudem geprüft statt
undefiniert (Abschnitt~\ref{asec:if-vector}). Für diese Arbeit ist das die zentrale Beobachtung:
Sämtliche Map-Zugriffe und -Modifikatoren der Hülle ruhen auf einem einzigen Such-Primitiv der
Modulgrenze --- der Punkt-Suche, flankiert von Einfügen und Löschen ---; eine Achsen-Implementierung
muss daher nur eine korrekte Punkt-Suche samt Knoten-Verkettung bereitstellen, der Rest ergibt sich
als feste Komposition. Die geordneten Operationen --- \texttt{begin}/\texttt{end},
\texttt{lower\_bound}, \texttt{upper\_bound}, \texttt{equal\_range} --- setzen zusätzlich eine
Ordnungs-Suche voraus und sind deshalb als eigene, additive Fähigkeit der Modulgrenze geführt;
ungeordnete Achsen erfüllen dieselbe Hülle ohne sie und bezahlen die fehlende Ordnung beim
geordneten Lesen mit linearem Aufwand. Genau diese Austauschbarkeit ist der zentrale Mess-Akt
(Abschnitt~\ref{sec:measurement-system}). Das Konformitäts-Gatter prüft je Gattung die in den
Tabellen~\ref{tab:if-map} und~\ref{tab:if-vector} mit $\dagger$ markierte Kern-Teilmenge.

\begin{longtable}{@{}>{\raggedright\arraybackslash}p{3.3cm} >{\raggedright\arraybackslash}p{5.6cm} >{\raggedright\arraybackslash}p{3.1cm}@{}}
  \caption{Zusammengesetzte Funktionen (\enquote{unter der Haube})}\label{tab:if-decomp}\\
  \toprule
  \textbf{Funktion} & \textbf{Zusammensetzung} & \textbf{Anmerkung} \\
  \midrule
\endfirsthead%
  \toprule
  \textbf{Funktion} & \textbf{Zusammensetzung} & \textbf{Anmerkung} \\
  \midrule
\endhead%
  \texttt{map::}\allowbreak\texttt{operator[]} & \texttt{lower\_bound} + Default-Einfügen bei Fehlen $\to$ Referenz & mutiert bei Fehlen; kein \texttt{out\_of\_range}, das Einfügen kann werfen \\
  \texttt{map::at} & \texttt{find} + Wurf \texttt{std::out\_of\_range} bei Fehlen & einziger Zugriff mit Wurf bei Fehlen; mutiert nie \\
  \texttt{map::insert} / \texttt{emplace} & Ordnungs-Suche + Knoten verlinken nur bei fehlendem Schlüssel & kein Überschreiben; \texttt{emplace} kann Argumente bei Treffer dennoch konsumieren \\
  \texttt{map::try\_\allowbreak{}emplace} & wie \texttt{emplace}, aber Argumente bei Treffer unangetastet & bewahrt move-only-Argumente \\
  \texttt{map::insert\_\allowbreak{}or\_\allowbreak{}assign} & Ordnungs-Suche + Einfügen oder Zuweisen & überschreibt (anders als \texttt{insert}) \\
  \texttt{map::count} / \texttt{contains} / \texttt{equal\_range} & reduzieren auf \texttt{find} bzw.\ \texttt{lower\_bound}+\texttt{upper\_bound} & bei Unikat-Map stets 0 oder 1 \\
  \texttt{vector::}\allowbreak\texttt{push\_back} / \texttt{emplace\_back} & ggf.\ Wachstum (Realloc + Umzug) + Konstruktion am Ende & versteckte Kosten + Invalidierung beim Realloc \\
  \texttt{vector::insert} / \texttt{emplace} & ggf.\ Wachstum + Schwanz verschieben + Einsetzen & zusätzlicher $O(n)$-Shift \\
  \texttt{vector::resize} & anhängen (Default/Wert) oder Schwanz kürzen & bidirektional; verkleinert die Kapazität nicht \\
  \texttt{vector::assign} & \texttt{clear} + Bereich/$n\times$ einfügen & Kapazität wird ggf.\ wiederverwendet \\
  \bottomrule
\end{longtable}

\noindent An der Modulgrenze ist die Primitivitäts-Ordnung dieser Tafel umgekehrt: Der eine
Einfüge-Antrieb der Map-Gattung trägt die Semantik von \texttt{insert\_or\_assign} --- Einfügen oder
Aktualisieren, die Rückgabe meldet, ob ein neuer Schlüssel entstand ---, sodass
\texttt{insert\_or\_assign} die Durchreichung ist und die nicht überschreibenden Formen
\texttt{insert}, \texttt{emplace} und \texttt{try\_emplace} host-seitig als Punkt-Suche mit
anschließendem Einfügen nur bei fehlendem Schlüssel komponiert werden: ein Pfad für drei
Standard-Funktionen, den das Gatter gegen drei Orakel-Aufrufe stellt. Der ISO-Vertrag führt es
umgekehrt (\texttt{insert} als Grundform, \texttt{insert\_or\_assign} als Zusammensetzung); beide
Sichten sind korrekt und beschreiben verschiedene Ebenen --- genau der Kontrast, der die
Austauschbarkeit hinter einer festen Schnittstelle konkret macht.

\noindent Funktionen, die beide Hüllen dem Namen nach \emph{teilen} --- \texttt{size}, \texttt{empty},
\texttt{clear}, die Iterator-Familien, \texttt{swap}, \texttt{get\_allocator} sowie
\texttt{insert}/\texttt{emplace}/\texttt{erase} ---, tragen dieselbe Rolle, aber verschiedene
Bezugsgrößen: in der Map Schlüssel-, im Vektor Positions-/Index-Semantik; \texttt{operator[]} etwa
fügt in der Map ein, ist im Vektor hingegen ein ungeprüfter Indexzugriff --- in der eigenen
Sequenz-Hülle geprüft (Abschnitt~\ref{asec:if-vector}), in der eigenen Map-Hülle host-seitig
komponiert (siehe oben). Diese Trennung ist in der Cache-Engine strukturell erzwungen: Map- und
Sequenz-Gattung besitzen getrennte Antriebs-Schnittstellen und getrennte Orakel, ein
gattungsfremder Adapter ist im Typsystem ausgeschlossen. Welche dieser Funktionen die Modulgrenze
trägt, welche der Host komponiert und welche nicht verdrahtet sind, verzeichnet das eingangs
genannte Entwickler-Dokument je Funktion mit ihrem Sprungpfad.
